\documentclass{article}
\usepackage[T1]{fontenc}
\usepackage{lmodern}
\usepackage{graphicx}
\usepackage{amsmath,amssymb}
\usepackage[a4paper,margin=2cm]{geometry}
\usepackage[absolute,overlay]{textpos}
\setlength{\TPHorizModule}{1mm}
\setlength{\TPVertModule}{1mm}
\usepackage[indonesian]{babel}
\usepackage{hyperref}
\setlength{\emergencystretch}{2em}
% unit: Halaman 23

\begin{document}

% === Halaman 23 (llm) ===

\begin{itemize}
    \item Selanjutnya,
    \begin{itemize}
        \item \texttt{Rtate} mungkin adalah \texttt{state},
        \item \texttt{atthattMZE} mungkin adalah \texttt{atthattime},
        \item \texttt{heVe} mungkin adalah \texttt{here}.
    \end{itemize}
    \item Jadi, kita memperoleh pemetaan baru:
    \[
    \begin{aligned}
    R & \rightarrow s \\
    M & \rightarrow i \\
    Z & \rightarrow m \\
    V & \rightarrow r
    \end{aligned}
    \]
\end{itemize}

\end{document}
